% !PRINTED 549
\weberpage{549}
We shall choose a positive quantity $\varrho_1$ and an angle $\Theta_1$ in such
a way that
\weq{9}{\varphi_1 = \varrho_1 e^{i\Theta_1}, \qquad \varphi_1' = \varrho_1 e^{-i\Theta_1},}
or
\weq{10}{\begin{aligned} \varphi_1 &= \varrho_1(\cos\Theta_1 + i\sin\Theta_1) \\ \varphi_1' &= \varrho_1(\cos\Theta_1 - i\sin\Theta_1), \end{aligned}}
from which it further follows that
\weq{11}{\varrho_1^2 = \varphi_1 \varphi_1'.}
Now
\weq{12}{(\varepsilon^{\lambda_1}, \alpha)(\varepsilon^{-\lambda_1}, \alpha) = \pm a_1}
is a quantity of the field $\Omega(\varepsilon)$ by (\S.\,164, 3.), and indeed,
since it does not change on passing to the conjugate imaginary value,
that is, on the interchange of $\varepsilon$ and $\varepsilon^{-1}$, it is a
real quantity. We shall determine the sign so that $a_1$ is
positive. Hence by (11), since $\varrho_1^2$ is positive,
\[
\varrho_1^2 = (\pm a_1)^m = a_1^m,
\]
and from this it follows that for odd $m$ the upper sign
holds in any case; for even $m$ the lower can also occur. We then have
\weq{13}{\varrho_1 = \sqrt{a_1^m},}
where the square root is to be taken positive.

Further, for the same reason,
\weq{14}{\frac{\varphi_1 + \varphi_1'}{2} = b_1, \qquad \frac{\varphi_1 - \varphi_1'}{2i} = c_1}
are real quantities in $\Omega(\varepsilon)$, and it follows from (10) that
\weq{15}{\cos\Theta_1 = \frac{b_1}{\sqrt{a_1^m}}, \qquad \sin\Theta_1 = \frac{c_1}{\sqrt{a_1^m}},}
from which the relation further follows:
\[
a_1^m = b_1^2 + c_1^2.
\]

Accordingly we obtain
\[
\sqrt[p_1]{\varphi_1} = \sqrt{a_1^{m_1}}\, e^{\frac{i\Theta_1}{p_1}};
\]
and one now proceeds in the same way with the functions $\varphi_2$ and so on.
One thus determines $\Theta_2 \ldots$ from a system of equations
like (15), and, if one further puts
\weq{16}{m_1\Theta_1 + m_2\Theta_2 + \cdots = \Theta}
one obtains, by (8)
%CONT>>
